% label-packages.tex -- P9: 日本語ラベルのパッケージ横断解決。
% 検証対象（plan ## Verification label行）:
%   * \ref/\pageref（star両形）/\eqref（amsmath経由）
%   * \autoref（hyperref）
%   * \nameref（nameref）
%   * \setcounterref/\getrefnumber（refcount; edef内完全展開）
%   * \cref/\cpageref（cleveref; TL2025に入っている）
%   * \hyperref[日本語ラベル]（リンク）
%   * 3パス収束（aux安定）
% 判定は tests/check-label-packages.py（抽出テキストの完全一致検査）。
\documentclass[pdflatex]{jlreq}
\usepackage{amsmath}
\usepackage[unicode]{hyperref}
\usepackage{refcount}
\usepackage{cleveref}
\pagestyle{empty}
\begin{document}
\section{日本語ラベルの検証}\label{日本語ラベル}
\label{sec:ascii}
\begin{equation}
e^{i\pi}+1=0 \label{式:オイラー}
\end{equation}
REF: \ref{日本語ラベル}|\ref*{日本語ラベル}|\pageref{日本語ラベル}|\pageref*{日本語ラベル}
EQREF: \eqref{式:オイラー}
AUTOREF: \autoref{日本語ラベル}
NAMEREF: \nameref{日本語ラベル}
CREF: \cref{日本語ラベル}|\cpageref{日本語ラベル}
HYPER: \hyperref[日本語ラベル]{リンク}
ASCII: \ref{sec:ascii}|\cref{sec:ascii}
\newcounter{rcnt}\setcounterref{rcnt}{日本語ラベル}
COUNTERREF: \arabic{rcnt}
GETREF: \getrefnumber{日本語ラベル}
\makeatletter
\edef\erefnum{\getrefnumber{日本語ラベル}}
\typeout{LABELPKG-EDEF: [\erefnum]}
\ifnum\erefnum=1 \typeout{LABELPKG-EDEF-OK}\else \typeout{LABELPKG-EDEF-FAIL}\fi
\makeatother
\end{document}
